\section{Random routing and independent tests}
\label{sec:routing}

We now use the windows from the preceding section to construct a random
periodic set \(B\) of expected density \(p\).
Selectors send each spatial point to one leaf of the tree, where an
independent terminal entry decides membership in \(B\).
For a stable center, the first default choice on its path will provide
several child windows in which to test nearby points.

\subsection{The tables and the selected set}

For each edge \(e=(P,i)\) with \(i<M\), let
\[
 S_e:\{0,\ldots,2^{b_e+2}-1\}\longrightarrow\{0,1\}
\]
be a random table whose entries are independent Bernoulli variables
with parameter \(1/2\).
There is no selector table for the default edge \((P,M)\).
For a leaf \(L\), let \(b(L)\) be the endpoint of the window assigned
to its unique incoming edge, and let
\[
 T_L:\{0,\ldots,2^{b(L)+2}-1\}\longrightarrow\{0,1\}
\]
be a random table whose entries are independent Bernoulli variables
with parameter \(p\).
All entries of all selector and terminal tables are mutually
independent.  The tree and every table are finite, so these variables
define a finite product probability space \(\Omega\).  We write
\(\mathbb P\) and \(\mathbb E\) for its probability and expectation.

Given the tables and a point \(z\in\mathbb R\), start at the root.
At an internal node \(P\), choose the first child \(P_i\), \(i<M\),
for which
\[
 S_{(P,i)}\bigl(J_{b_{(P,i)}}(z)\bigr)=1.
\]
If all \(M-1\) values are zero, choose the default child \(P_M\).
Continue until reaching a leaf, denoted by \(L(z)\), and define
\begin{equation}
 \label{eq:random-set}
 B=\left\{z\in\mathbb R:
 T_{L(z)}\bigl(J_{b(L(z))}(z)\bigr)=1\right\}.
\end{equation}
The table rules use only periodic keys.  Since all grids in the
construction are nested, membership in \(B\) is constant on each cell
of the finest grid used.  Thus every realization of \(B\) is a
\(1\)-periodic measurable set, consisting of finitely many half-open
cells per period.

Let \(\mathcal S\) be the sigma-field generated by every entry of every
selector table.  For fixed \(z\), conditioning on \(\mathcal S\) fixes
\(L(z)\) and the key of its terminal entry.  That entry still has
probability \(p\) of being one, since terminal tables are independent
of all selectors.  Hence
\[
 \mathbb P(z\in B\mid\mathcal S)=p.
\]
Integrating over one period and interchanging the integral with the
finite weighted sum over \(\Omega\) gives
\begin{equation}
 \label{eq:mean-density}
 \mathbb E\rho(B)=\int_0^1\mathbb P(z\in B)\,dz=p.
\end{equation}

\subsection{Conditioning at a center}

For a fixed center \(x\in\mathbb R\), expose one entry from every
selector table, including those in subtrees not visited by its path.
The resulting sigma-field is
\begin{equation}
 \label{eq:center-exposure}
 \mathcal F_x
 =
 \sigma\left(
 S_e\bigl(J_{b_e}(x)\bigr):
 e=(P,i),\ i<M
 \right)
 \subseteq\mathcal S.
\end{equation}
An atom \(\xi\) of \(\mathcal F_x\) means the event specifying all
these exposed values; we write \(\mathbb P(\,\cdot\mid\xi)\) for
conditioning on that event.
Every such atom has positive probability.
The exposed addresses \((e,J_{b_e}(x))\) are deterministic once
\(x\) is fixed, with the table name \(e\) included in the address.
Consequently, on each atom, all remaining selector entries retain
their independent fair law, and all terminal entries retain their
independent Bernoulli-\(p\) law.

The exposed values determine the full path of \(x\).
When that path chooses a default child, let \(U\) denote its first
internal node at which this happens.  Thus the existence and identity
of \(U\) are determined by \(\mathcal F_x\).
The next estimate concerns the original product law, before an atom
of \(\mathcal F_x\) is fixed.

\begin{lemma}[A default choice on the center path]
 \label{lem:first-default}
For every fixed \(x\in\mathbb R\),
\[
 \mathbb P\bigl(\text{the path of \(x\) never chooses child \(M\)}\bigr)
 =
 (1-2^{1-M})^d<p.
\]
\end{lemma}

\begin{proof}
At any internal node, the default child is chosen exactly when its
\(M-1\) selector entries at \(x\) are all zero.  This has probability
\(2^{-(M-1)}=2^{1-M}\).

To compute the probability along the path, reveal the entries at an
internal node only when the path first reaches it.  The decisions
leading to that node use tables at its strict ancestors.  Its own
tables are distinct from all of those tables and have not yet been
revealed, so their entries at \(x\) remain independent fair bits,
conditional on the preceding path history.
At each of the \(d\) internal-node decisions, the conditional
probability of a non-default choice is therefore \(1-2^{1-M}\).
Iterating this conditional probability gives
\((1-2^{1-M})^d\).  The strict inequality is the choice of \(d\) in
\eqref{eq:tree-parameters}.
\end{proof}

\subsection{Routing trials through the first default}

Fix henceforth a stable center \(x\in G\), and condition on an atom
\(\xi\) of \(\mathcal F_x\) for which \(U\) exists.
The node \(U\) is fixed under this conditioning.  Let \(h\) be its
height, and write
\[
 r=r_h,\qquad e_i=(U,i),\qquad b_i^*=b_{e_i}^*
 \quad(1\le i<M),\qquad y_n(t)=x+t2^{-n}.
\]
At \(U\), all the exposed non-default selector values are zero.

For \(i<M\) and \(z\in\mathbb R\), let \(L_i(z)\) be the leaf obtained
by applying the same path rule starting at the child \(U_i\).
If \(U_i\) is already a leaf, set \(L_i(z)=U_i\).
Define the local predicate
\begin{equation}
 \label{eq:local-predicate}
 Q_i(z)=
 S_{e_i}\bigl(J_{b_{e_i}}(z)\bigr)\,
 T_{L_i(z)}\bigl(J_{b(L_i(z))}(z)\bigr).
\end{equation}
The function \(Q_i\) takes values in \(\{0,1\}\) and consults only
the selector on \(e_i\) and tables in the subtree rooted at \(U_i\).
It provides a sufficient condition for membership in \(B\) at the
points indexed by \(W_{e_i}\).

\begin{lemma}[Preservation of the route]
 \label{lem:routing}
For every outcome in \(\xi\), every \(i<M\), every \(n\in W_{e_i}\),
and every \(t\in[1,2]\), the actual root path of \(y_n(t)\) reaches
\(U\) and rejects children \(1,\ldots,i-1\) there.
In particular,
\begin{equation}
 \label{eq:local-hit}
 Q_i(y_n(t))=1
 \quad\Longrightarrow\quad y_n(t)\in B.
\end{equation}
\end{lemma}

\begin{proof}
Consider a strict ancestor \(P\) of \(U\), and let \(P_a\) be its
child on the path of \(x\) to \(U\).
Since \(U\) is the first default node, \(a<M\).
The decision at \(P\) uses precisely the zero values of the selectors
\(1,\ldots,a-1\) and the value one of selector \(a\).
Each of the edges \((P,j)\), \(j\le a\), precedes every edge in the
subtree of \(P_a\) in the prescribed preorder.
Their windows therefore precede \(W_{e_i}\).
By the stable-key identity \eqref{eq:stable-keys},
each of these selector evaluations at \(y_n(t)\) agrees with its
evaluation at \(x\).
Induction from the root along the ancestor chain shows that
\(y_n(t)\) makes the same choices and reaches \(U\).
Later sibling selectors at an ancestor are not needed to make its
choice.

For every \(j<i\), the edge \((U,j)\) also precedes \(e_i\).
Its selector value at \(x\) is zero, because \(x\) takes the default
child at \(U\).
Another application of \eqref{eq:stable-keys} makes the corresponding
value at \(y_n(t)\) zero.
Thus the actual path rejects children \(1,\ldots,i-1\) at \(U\).

If \(Q_i(y_n(t))=1\), the selector on \(e_i\) is one, so the actual
path chooses child \(i\) at \(U\).
Its continuation is exactly the rule defining \(L_i(y_n(t))\),
and the terminal entry in \eqref{eq:local-predicate} is one.
The membership rule \eqref{eq:random-set} now gives \(y_n(t)\in B\).
Figure~\ref{fig:routing} records this common-prefix and branching argument.
All equalities used here hold for every completion of the unexposed
tables and every \(t\in[1,2]\), as required.
\end{proof}


\begin{figure}[!htb]
\centering
\begin{tikzpicture}[x=1cm,y=1cm,font=\small,>=Latex]
 \node[draw,rounded corners,fill=gray!8,minimum width=4.3cm] (prefix) at (0,3.15)
   {common choices at strict ancestors};
 \node[draw,circle,inner sep=3pt] (u) at (0,1.85) {$U$};
 \draw[->,thick] (prefix)--node[right] {stable earlier keys}(u);
 \node[draw,rounded corners,fill=blue!8] (childi) at (-3,.15) {$U_i$, $i<M$};
 \node[draw,rounded corners,fill=orange!12] (childm) at (3,.15) {$U_M$};
 \draw[->,thick,blue!70!black] (u)--node[above left,align=right]
   {$y_n(t)$: selector $i$ succeeds\\earlier children rejected}(childi);
 \draw[->,thick,orange!80!black] (u)--node[above right,align=left]
   {$x$: all $M-1$ selectors vanish\\default choice}(childm);
 \node[align=center,below=5mm of childi] (terminal)
   {local route to $L_i(y_n(t))$\\terminal entry equals $1$};
 \draw[->,blue!70!black] (childi)--(terminal);
\end{tikzpicture}
\caption{Route preservation for $x\in G$, $n\in W_{(U,i)}$ and
$t\in[1,2]$, conditional on $Q_i(y_n(t))=1$.
The center $x$ and its translate
$y_n(t)$ make the same required choices above the first default node $U$.
At $U$, the center takes child $M$, while the translate rejects the earlier
children and takes child $i$. The local terminal success then places
$y_n(t)$ in $B$. The diagram is conditional on the successful local
predicate; it does not assume independence of the chosen leaves.}
\label{fig:routing}
\end{figure}

